% Generated draft: requires mathematical and source review.
\documentclass[11pt]{article}
\usepackage[margin=25mm]{geometry}
\usepackage{fontspec}
\setmainfont{Latin Modern Roman}
\usepackage{amsmath,amssymb,mathtools,unicode-math}
\setmathfont{Latin Modern Math}
\usepackage{enumitem,needspace,xurl,hyperref}
\setlength{\parindent}{0pt}
\setlength{\parskip}{5pt}
\setlength{\emergencystretch}{3em}
\newcommand{\N}{\mathbb{N}}
\newcommand{\Z}{\mathbb{Z}}
\newcommand{\Q}{\mathbb{Q}}
\newcommand{\R}{\mathbb{R}}
\newcommand{\C}{\mathbb{C}}
\newcommand{\F}{\mathbb{F}}
\newcommand{\A}{\mathbb{A}}
\newcommand{\PP}{\mathbb{P}}
\newcommand{\E}{\mathbb{E}}
\newcommand{\eps}{\varepsilon}
\newcommand{\dd}{\,\mathrm{d}}
\begin{document}
\section*{UnsolvedMath ID: 30005196}
\textbf{Source problem code:} catalogue entry
\subsection{Definitions and mathematical statement}
The supplied entry poses four related research requests rather than one fully specified conjecture. The following definitions give conventional meanings to its principal terms; the qualifications concerning explicitness, algebraicity, and analogous results remain unresolved.

Let \(\operatorname{RE}\) denote the class of languages recognized by Turing machines that halt and accept exactly on members of the language. Under the usual conventions, \(\operatorname{MIP}^{*}\) denotes the class of languages admitting interactive proofs with a probabilistic polynomial-time classical verifier and multiple noncommunicating provers sharing prior quantum entanglement. Strategies may use arbitrarily large finite-dimensional entangled states. Completeness and soundness have a fixed positive gap, conventionally at least \(2/3\) versus at most \(1/3\). The record names the equality \(\operatorname{MIP}^{*}=\operatorname{RE}\) as the result whose proof is to be refined; it does not specify a particular protocol model or complexity-theoretic normalization.

Let \(R\) be the hyperfinite type \(\mathrm{II}_1\) factor with normalized trace \(\tau_R\), and let \(\omega\) be a free ultrafilter on \(\N\). Its tracial ultrapower is
\[
 R^{\omega}=\ell^{\infty}(\N,R)\big/\{(x_k):\lim_{k\to\omega}\tau_R(x_k^*x_k)=0\},
 \qquad
 \tau_{\omega}([(x_k)])=\lim_{k\to\omega}\tau_R(x_k).
\]
A standard formulation of the Connes embedding problem asks whether every finite von Neumann algebra \(M\) with separable predual and a faithful normal tracial state \(\tau\) admits a unital injective \(*\)-homomorphism \(\iota:M\to R^{\omega}\) satisfying \(\tau_{\omega}\circ\iota=\tau\). Thus a standard precise form of a requested counterexample is a specified pair \((M,\tau)\) satisfying these hypotheses for which no such embedding exists. If the intended output is instead a unital separable \(C^*\)-algebra \(A\) with a tracial state \(\tau\), the corresponding requirement is that its tracial GNS von Neumann algebra \(\pi_{\tau}(A)''\), with the induced trace, fail this embedding property. These are possible formulations of the request; the supplied entry does not select the presentation category or define how explicit the algebra and trace must be.

For a discrete group \(\Gamma\) with identity \(e\), write \(\operatorname{tr}_n=n^{-1}\operatorname{Tr}\) and \(\|T\|_{2,n}=\operatorname{tr}_n(T^*T)^{1/2}\). The group is hyperlinear if, for every finite \(F\subseteq\Gamma\) containing \(e\) and every \(\varepsilon>0\), there exist \(n\geq1\) and a map \(u:F\to U(n)\) such that
\[
 u(e)=I_n,\qquad
 \|u(gh)-u(g)u(h)\|_{2,n}<\varepsilon
 \quad(g,h,gh\in F),
\]
\[
 |\operatorname{tr}_n(u(g))|<\varepsilon
 \quad(g\in F\setminus\{e\}).
\]
It is nonhyperlinear if this approximation property fails, equivalently if some finite \(F\) and some \(\varepsilon>0\) admit no such \(n\) and \(u\). This definition does not impose countability, finite generation, or finite presentation; the record supplies none of those restrictions.

\begin{enumerate}
\item Give a more algebraic proof of \(\operatorname{MIP}^{*}=\operatorname{RE}\). The source supplies no criterion for comparing the algebraicity of proofs.
\item Construct an explicit counterexample to the Connes embedding problem, including a specification of its tracial state. The intended category and standard of explicitness require clarification as described above.
\item Identify analogous results or consequences. The source does not specify the objects, hypotheses, or conclusions to which the analogy should apply, so this part cannot be stated as a definite quantified assertion.
\item Construct a nonhyperlinear group. The source does not prescribe a presentation or a standard of explicit construction.
\end{enumerate}

For context, the original excerpt also contains a corrupted nonlocal-game formula. Under the conventional finite two-player interpretation, with question sets \(X,Y\), answer sets \(A,B\), question distribution \(\pi\), acceptance predicate \(V:A\times B\times X\times Y\to\{0,1\}\), and conditional probabilities \(p(a,b\mid x,y)\geq0\) summing to one over \(a,b\), the acceptance probability is
\[
 \omega(G,p)=\sum_{x\in X,\,y\in Y}\sum_{a\in A,\,b\in B}
 \pi(x,y)V(a,b\mid x,y)p(a,b\mid x,y).
\]
This is a conventional reconstruction of the contextual formula. The damaged excerpt does not reliably specify its alphabets, the meaning of its \(\alpha\) subscripts, or the admissible strategy class.
\subsection{Short English statement}
The entry asks for a more algebraic proof of \(\operatorname{MIP}^{*}=\operatorname{RE}\), an explicit Connes embedding counterexample with its trace, related results, and a construction of a nonhyperlinear group. Several requests remain qualitative: the source does not define algebraicity, explicitness, or the intended analogies. Its partially solved label and literature assessment are supplied dated curation, not independently verified conclusions.
\subsection{Sources}
\begin{itemize}
\item Supplied local UnsolvedMath registry record, permanent ID \(30005196\), source problem code \texttt{catalogue entry}, in the Oberwolfach Reports Open Problems collection. Local registry locations identified by the project are \texttt{data/unsolvedmath/problems.json} and the supplemental \texttt{catalog/problems.json}. No individual proposer is named in the supplied record.
\item Supplied report citation: \(C^*\)-Algebras (2023), Oberwolfach Reports, DOI \texttt{10.4171/owr/2022/36}, pp.\ 2059--2127. Original report link: \url{https://doi.org/10.4171/owr/2022/36}. The report label gives 2023, whereas the DOI and supplied proposed year point to 2022; this bibliographic discrepancy is unresolved here.
\item Supplied literature source: \emph{MIP-star equals RE}, \url{https://arxiv.org/abs/2001.04383}.
\item Supplied literature source: \emph{The Connes Embedding Problem, a guided tour}, \url{https://arxiv.org/abs/2109.12682}.
\item Literature assessment dated exactly \texttt{2026-08-22}: the supplied curation labels the entry partially solved and describes \(\operatorname{MIP}^{*}=\operatorname{RE}\) as giving an explicit negative resolution of the Connes embedding problem, while describing a concrete nonhyperlinear group and the other algebraic refinements as remaining open. This is supplied dated curation, not an independently verified literature assessment. The background expressly requests independent verification of time-sensitive status and cited claims.
\item The registry labels its statement \texttt{corrected\_verified} and reports that it was reconstructed and checked against the official Oberwolfach report. That verification claim is supplied metadata; no source was visited or independently checked in preparing this draft.
\end{itemize}
\end{document}
